%!TEX program = lualatex
\documentclass[a5paper,ngerman,headheight=80pt,10pt,DIV=calc]{scrartcl}
\usepackage[ngerman]{babel}
\usepackage{blindtext}
\usepackage[utf8]{luainputenc}
\usepackage{scrlayer-scrpage}
\usepackage[T1]{fontenc}
\usepackage{fontspec}
\usepackage{multicol}
\usepackage{enumitem}
\usepackage{epigraph}
\usepackage[german=quotes]{csquotes}
\usepackage[footskip=20pt,bottom=50pt]{geometry}
\setmainfont{Source Sans Pro}
\MakeOuterQuote{"}

\setlength{\parindent}{0pt}
\setlength{\parskip}{4pt}

\pagestyle{plain.scrheadings}
\thispagestyle{scrheadings}
\lohead{\vspace*{-1.1cm}\includegraphics[scale=0.3]{../images/stupa_logo_rot.pdf}}
\rohead{\\Präsidium des\\Studierendenparlaments}

\lofoot[]{V.i.S.d.P. Gunhild Berdal, Jim Martens, Ramon Weilinger\\Von-Melle-Park 5, 20146 Hamburg}
\cofoot[]{}
\cohead[\thepage]{}

\renewcommand{\thesubsection}{TOP \arabic{subsection}}
\renewcommand{\thesubsubsection}{\arabic{subsubsection}.}
\setcounter{subsection}{-1}

\begin{document}
    UHH · StuPa-Präsidium · Von-Melle-Park 5 · D-20146 Hamburg\\
    Tel.: 040/450 204 39 · stupa@uni-hamburg.de

    \subsection*{Infos zur Wahl zum Studierendenparlament 2017/18}

    Nächste Woche bekommst Du per Post die Wahlunterlagen für die Briefwahl
    zum Studierendenparlament 2017/18 zugeschickt. In dem Briefumschlag befinden
    sich ein Anschreiben mit Wahlschein, eine
    Wahlbroschüre mit Informationen zur Wahl und den Darstellungen der Listen,
    ein gelber Stimmzettelumschlag, ein weißer Rückumschlag und ein gelber
    Stimmzettel. Um eine gültige Stimme abzugeben, beachte bitte die Hinweise
    in der Wahlbroschüre. Spätestens am 2. Januar 2017 muss der Rückumschlag
    bei uns angekommen sein.

    Vom 16. bis 20. Januar 2017 besteht die Möglichkeit auch persönlich an einer
    Urne die Stimme für das Studierendenparlament abzugeben. Es kann insgesamt
    nur \textbf{EINE} Stimme abgegeben werden -- entweder per Brief oder per
    Urne.

    Im Verlaufe des Dezembers wirst Du auch die Wahlunterlagen für die Wahl
    zum Akademischen Senat bekommen. Die Wahl zum Akademischen Senat findet
    als reine Briefwahl statt. Bitte achte auf die Inhalte und wirf nicht
    versehentlich Wahlunterlagen -- und Deine Stimme -- weg.

    Wenn Deine Eltern die Post bekommen, sage ihnen bitte Bescheid, sodass Du
    die Wahlunterlagen bekommst und von Deinem Wahlrecht Gebrauch machen kannst.

    Damit Du eine informierte Entscheidung treffen kannst, gibt es am 15. Dezember
    eine Listenvorstellung, auf der die antretenden Listen zu grundsätzlichen
    weltanschaulichen Fragen (z.B. "Wie soll Hamburg in zehn Jahren aussehen?",
    "Wie soll die Gesellschaft aussehen?" und "Wie soll Hochschule aussehen?")
    Rede und Antwort stehen.

    \begin{center}
        \textbf{
        Erste Listenvorstellung\\
        Donnerstag, 15. Dezember 2016,\\
        18 Uhr\\
        im Raum 0079 des WiWi-Bunkers\\
        (Von-Melle-Park 5, 20146 Hamburg)}
    \end{center}
\end{document}
